\section{Background: Recursive Spawning}
\label{sec:background-recursive-spawning}

The Bailout Protocol — defined formally in Section~\ref{sec:bailout} — requires a substrate in which every agent node is reachable from a human principal through a defined, immutable delegation chain. That substrate is the \bxthree{} Framework. This section establishes the four enabling lineages that make recursive spawning of agent chains tractable: agent lifecycle management in multi-agent systems, accountability and provenance in distributed AI deployments, the distinction between fixed and mutable delegation architectures, and hierarchical task decomposition as the operative paradigm for planning across spawned agent hierarchies. Together these motivate the \bxthree{} immutable parent-pointer architecture as the substrate on which the Bailout Protocol operates.

% ---------------------------------------------------------------
\subsection{Agent Lifecycle Management in Multi-Agent Systems}
\label{sec:lifecycle}

An \emph{agent lifecycle} comprises the states an agent occupies from creation through termination, the events that drive transitions between those states, and the resources and permissions the agent holds at each stage. In single-agent systems the lifecycle is straightforward: a principal instantiates an agent, grants it a fixed capability set, and retires it when its task is complete. In multi-agent systems the lifecycle becomes recursive: an agent $A$ may spawn a sub-agent $B$ to handle a subtask, granting $B$ a subset of its own capabilities. When $B$ in turn spawns $C$, the delegation chain deepens, and the lifecycle of $C$ is nested inside that of $B$, which is nested inside that of $A$.

This recursive spawning is described in the agent networks literature as a \emph{delegation chain} \cite{ietf-irtf-agent-delegation}. The IETF IRTF note on AI agent communication identifies four structural challenges that compound with each additional hop in such chains: scope attenuation (each delegation should narrow permissions, but the mechanism for enforcing this at runtime is frequently absent), transitive accountability (the initiating user or system must remain reachable even after $N$ hops of delegation), provenance continuity (every action taken by every spawned agent must be traceable back to an authorized parent), and identity interoperability (different agent frameworks use incompatible identity systems, making cross-framework delegation chains fragile or unverifiable) \cite{ietf-irtf-agent-delegation}. These challenges are not hypothetical: production multi-agent deployments routinely encounter them when sub-agents spawn further sub-agents to handle exception conditions or to decompose temporally extended tasks \cite{csa-delegation, dynamiq-multiagent}.

Conventional agent lifecycle management frameworks treat spawning as a dynamic, runtime decision made by the parent agent, with the runtime responsible for tracking which agents exist, which are active, and which hold which permissions. This dynamic model is powerful but introduces a bootstrapping problem for accountability: if the parent chain is mutable — if agents can redirect their parent pointers, revoke sub-agent permissions, or terminate sub-agents without a defined protocol — then the accountability chain is also mutable, and any audit conducted after the fact may find that the chain of delegation has been altered, compressed, or deleted. Static or append-only lifecycle management avoids this problem at the cost of flexibility; dynamic lifecycle management enables rich multi-agent behavior at the cost of provenance integrity unless constrained by an architectural layer that enforces immutability of the delegation record \cite{compel-audit, agility-accountability}.

The \bxthree{} Framework resolves this tension by separating agent lifecycle management (which remains dynamic and capability-driven) from delegation record-keeping (which is append-only and Fact Layer-enforced). An agent may spawn a sub-agent at any time; the spawning event is immediately recorded in the Forensic Ledger with the parent identifier, the sub-agent identifier, the granted capability scope, and a timestamp. The record is immutable after write. This means that even if the parent agent is later compromised, retired, or refactored, the historical delegation chain is preserved intact and auditable. The Bailout Protocol depends on this separation: it traverses the immutable delegation record — not the runtime process table — to propagate exceptions upward toward the human anchor.

% ---------------------------------------------------------------
\subsection{Accountability and Provenance in Distributed Systems}
\label{sec:provenance}

Accountability in distributed systems is the property that every action can be attributed to an actor, and every actor can be attributed to a principal who bears responsibility for the outcome. Provenance is the complementary property that the complete causal history of a data item or decision — every transformation, aggregation, and derivation — can be reconstructed from an immutable record. In conventional distributed systems (databases, workflows, microservices), accountability and provenance are typically implemented through audit logs: append-only records of events, transactions, and state changes that together encode the causal chain \cite{compel-audit, augment-audit}.

Multi-agent AI systems amplify both properties to a degree that conventional audit log architectures cannot accommodate. In a microservices architecture, the causal chain of a decision is a sequence of service calls, each attributable to a principal via authentication credentials. The chain is linear (or tree-structured, in the case of fan-out calls) and its elements are deterministic: each service call is a discrete, enumerable action with a defined input and output. In a multi-agent system, the causal chain of a decision may traverse LLM inference calls, tool invocations, inter-agent message exchanges, spawned sub-agents, and human inputs — many of which are non-deterministic, context-dependent, and not reducible to a discrete service call with a clean input-output boundary \cite{arxiv-audit-trails, swept-audit}.

The enterprise AI audit literature identifies two criteria as necessary (though not sufficient) for a multi-agent output to pass regulatory or contractual audit: \emph{attributability} — every output segment traces to a specific agent, model version, and authorizing specification — and \emph{reversibility} — every output can be rolled back to a known-good prior state without cascading failure \cite{augment-audit}. Attributability requires that the delegation chain from the human principal to the acting agent be fully populated and immutable; reversibility requires that the state machine governing agent transitions have well-defined rollback semantics. Both criteria map directly to \bxthree{} architectural commitments: the immutable parent pointer ensures attributability, and the Bailout Protocol's state machine ensures that exception conditions have a defined resolution pathway rather than an undefined cascade.

The COMPEL framework for decision provenance in multi-agent systems establishes that flat event logs — sufficient for single-agent sequential reasoning traces — are inadequate for multi-agent architectures where concurrent agents interact, spawn sub-agents, and operate under partial observability \cite{compel-audit}. The framework proposes purpose-built provenance graphs that capture causal relationships between agents, not merely temporal ordering of events. The \bxthree{} Forensic Ledger is architecturally aligned with this requirement: it records causal relationships (parent-child spawn relationships, delegation scope, exception propagation paths) as first-class record fields, not merely as inferred ordering from timestamps.

Blockchain-based approaches to AI governance and attribution have been proposed as a mechanism for achieving immutability and provenance integrity in AI decision systems \cite{flexblok-blockchain, medium-accountability}. The core insight — that append-only, cryptographically sealed records are necessary for auditability in adversarial or semi-autonomous settings — is sound and aligns with the \bxthree{} Forensic Ledger's SHA-256 sealing mechanism. The \bxthree{} Framework differs architecturally in its placement of the immutable record: blockchain-based approaches typically add an external immutable layer atop existing systems; the \bxthree{} Framework makes the immutable record a native component of the agent state machine, co-located with the agent's execution context, and enforced by the \fact{} layer rather than appended by a separate subsystem.

% ---------------------------------------------------------------
\subsection{Fixed vs. Mutable Delegation Chains}
\label{sec:delegation-chains}

A \emph{delegation chain} is a sequence of nodes $(n_0, n_1, \dots, n_k)$ in which $n_i$ delegates authority to $n_{i+1}$ for some subset of its provisioned capabilities. The chain terminates at $n_k$, whose terminal capability set defines the actual operating authority at the leaf of the delegation. In human organizations, delegation chains are typically mutable: a manager may re-delegate authority, revoke a sub-ordinate's access, or reorganize reporting relationships. In software access-control systems, delegation chains are often fixed at provisioning time and enforced by reference monitors that check every access request against the statically defined permission set \cite{csa-delegation}.

The distinction between fixed and mutable delegation has significant consequences for the accountability properties of the resulting system. A fixed delegation chain is easier to audit (the permission graph is static and fully enumerated), but it is less adaptive: if the lead agent fails or the task requires a capability not pre-provisioned in the static grant, the chain cannot self-repair. A mutable delegation chain is more adaptive, but each mutation is a potential point of accountability loss: if the chain can be redirected after the fact, an auditor cannot trust that the delegation record reflects the actual authority state at the time of an action.

The AI agent delegation literature identifies mutable delegation chains as a primary attack surface. The Cloud Security Alliance's analysis of AI agent delegation identifies \emph{agent session smuggling} — in which a malicious actor exploits the gap between a delegate's initial authorization and its current operational context to expand the delegate's effective authority beyond the scope of the original grant — as a structural vulnerability in mutable delegation models \cite{csa-delegation}. The analysis recommends scope attenuation (each hop narrows permissions), cryptographic chain integrity (each delegation hop is cryptographically bound to the previous), and offline attenuation verification (downstream agents can verify the chain without consulting the originating system). These three properties are precisely what an immutable parent-pointer architecture provides architecturally.

The \bxthree{} Framework implements a \emph{fixed delegation record with dynamic capability grants}: the parent pointer $\alpha(n)$ is immutable after the spawn event is recorded, but the capability scope granted at each hop can be narrower than the parent's scope and can be further restricted at any subsequent re-delegation event. This preserves adaptability in capability management while ensuring that the delegation topology — who spawned whom, and what scope was granted at each hop — is frozen at the moment of each spawn and cannot be retroactively altered. The Bailout Protocol traverses the fixed topology; the dynamic capability layer governs what each node may do. These are independent dimensions of the architecture, and keeping them separate avoids the rigidity of purely static systems without incurring the accountability risks of purely dynamic ones.

DeepMind's work on intelligent AI delegation distinguishes between chains in which accountability is transitive and provenance is immutable — where each hop is cryptographically sealed and the chain can be reconstructed at any future point — and chains in which these properties do not hold \cite{deepmind-delegation}. The \bxthree{} immutable parent-pointer architecture is explicitly designed to realize the former property: the chain of delegation is append-only, each spawn event is a Fact Layer record with SHA-256 seal, and the chain is traversable by the Bailout Protocol at any point in the agent's lifecycle. This immutability is not a governance convention but an architectural enforcement: no machine actor, regardless of capability level, can alter a recorded spawn event without creating a ledger entry that is itself immutable and auditable.

% ---------------------------------------------------------------
\subsection{Hierarchical Task Decomposition in AI}
\label{sec:htn}

Hierarchical Task Decomposition (HTD) is the paradigm in which a complex goal is recursively decomposed into simpler sub-goals until the sub-goals are directly executable by primitive agents or actions. HTD is the dominant planning paradigm in classical AI \cite{erol1994umcp, georgievski2015htn} and has been re-adopted as the operative planning model for modern LLM-based agent systems, where it addresses the fundamental problem that a flat reasoning trace cannot efficiently represent the branching, parallel, and conditionally dependent structure of a multi-step complex task \cite{zylos-htn, agentorchestra, reagentree}.

Hierarchical Task Network (HTN) planning formalizes HTD by structuring the planning problem as a graph of tasks, where compound tasks decompose into subtask networks via methods, and primitive tasks execute directly as actions \cite{erol1994umcp, nau-htn-overview}. The planner begins with a high-level task — the root goal — and recursively expands compound tasks by selecting applicable methods, each of which specifies a partial ordering of subtasks. The process terminates when all tasks in the network are primitive. HTN planning is strictly more expressive than classical STRIPS-style planning under standard definitions of expressivity: certain goal-class tasks, such as constructing a complex artifact or planning a multi-leg journey, cannot be compactly represented as a goal state to be achieved but can be naturally represented as a task hierarchy to be decomposed \cite{erol1994semantics}.

Modern LLM agent research has rediscovered HTN as a structuring principle for multi-agent coordination. The ReAcTree framework demonstrates that decomposing complex tasks into a tree of LLM-powered agent nodes — where each node handles a subgoal and may recursively decompose further — outperforms flat reasoning on tasks with conditional branching, parallel subtask execution, and modular failure recovery \cite{reagentree}. AgentOrchestra formalizes a hierarchical multi-agent architecture in which a Planning Agent decomposes high-level directives and distributes sub-plans to Specialist Agents, each of which executes its sub-plan and reports results to the planner for synthesis and potential replanning \cite{agentorchestra}. The Zylos Research analysis of goal decomposition and hierarchical planning identifies four complementary techniques that have converged in contemporary agent systems: decomposition strategies (how goals are split), hierarchical planning structures (how subtasks are organized), execution patterns (how plans interact with real-world feedback), and replanning mechanisms (how plans are repaired when reality diverges from assumptions) \cite{zylos-htn}.

Recursive Task Decomposition (RTD) — the specific pattern in which an agent decomposes a goal into sub-agents, each of which may further decompose its sub-goal into further sub-agents — is the operative pattern for spawning in multi-agent AI systems \cite{github-ai-agent-design}. RTD enables parallel subtask execution (multiple sub-agents operate concurrently), localized failure recovery (a sub-agent failure is contained within its own subtree and does not cascade to sibling sub-agents), and modular reusability (the same sub-agent class can be instantiated across different parent tasks). These properties are the architectural rationale for why spawning — rather than sequential planning — is the preferred scaling strategy for complex, temporally extended tasks in multi-agent deployments.

The critical connection between RTD and accountability is the parent chain: when a sub-agent spawns a further sub-agent, the grandparent chain must remain intact for the Bailout Protocol to function. If sub-agent $B$, spawned by $A$, spawns sub-agent $C$, and $C$ encounters an exception, the exception must propagate to $A$ (or to $h(A)$, the human anchor of $A$), not merely to $B$, because $B$'s own resolution authority may be insufficient to handle the exception. This requires that the delegation chain be visible to the protocol at all depths, and that each node's human anchor be reachable through the chain. The \bxthree{} immutable parent-pointer architecture provides exactly this: every spawned node has a defined $\alpha(n)$ pointing to its parent, and every node carries $h(n)$ identifying its human anchor, enabling the Bailout Protocol to traverse arbitrary delegation depth without losing the accountability thread.

% ---------------------------------------------------------------
\subsection{The BX3 Framework as Substrate for Immutable Parent Chains}
\label{sec:bx3-immutable-chains}

The foregoing four lineages — lifecycle management, accountability and provenance, delegation chain immutability, and hierarchical task decomposition — converge in the \bxthree{} Framework's treatment of the agent graph as an immutable-parent-chain data structure. The \bxthree{} three-layer model (Purpose, Bounds, Fact) was described in \S\ref{sec:three-layer-model}; here we focus on the specific architectural properties of the parent chain that make it suitable as a bailout propagation substrate.

The \fact{} layer maintains the agent graph as a set of append-only records. Each spawn event creates a new record containing: the child node identifier $n$, the parent node identifier $\alpha(n)$, the granted capability scope $c(n)$, the human anchor identifier $h(n)$, and a SHA-256 hash of the record content sealed at write time. The record is immutable: no machine actor, including the parent node $\alpha(n)$, can modify or delete the record after write. The hash chain ensures that any tampering with a historical record is detectable: the hash of record $i$ depends on the hash of record $i-1$, forming a provably unbroken chain from the current record back to the genesis record.

This immutability property is critical for the Bailout Protocol's correctness. The protocol traverses the parent chain $(n, \alpha(n), \alpha(\alpha(n)), \dots)$ to propagate exceptions upward. If any link in this chain were mutable — if a compromised or faulty node could redirect its parent pointer, delete its spawn record, or insert itself between two existing nodes — the protocol could propagate an exception to a wrong node, fail to reach the human anchor, or lose the exception entirely. The \fact{} layer's append-only enforcement ensures that none of these failures can occur: the chain as it existed at spawn time is the chain that exists at any future point, and the Bailout Protocol traverses a topology that cannot be altered post-spawn.

The \bxthree{} spawn operation is formally defined as:

\begin{align}
\text{Spawn}(p, S, c, h) &\rightarrow n \\
\text{subject to:} \quad \alpha(n) &= p \\
\text{record}(n) &= \{p, S, c, h, t, \text{hash}(\text{record}(n))\} \\
\text{hash}(\text{record}(n)) &= \text{SHA256}(p \parallel S \parallel c \parallel h \parallel t \parallel \text{hash}(\text{record}(p)))
\end{align}

where $p$ is the parent node, $S$ is the sub-agent specification, $c$ is the granted capability scope, $h$ is the human anchor, $t$ is the spawn timestamp, and the record's hash is computed as a chain hash that includes the parent's record hash, ensuring that the spawn is causally linked to the entire upstream chain.

This definition makes the spawn operation \emph{non-destructive of provenance}: the parent chain is never modified by the spawn; a new node is appended with a reference to its parent. The chain grows monotonically. The Bailout Protocol therefore operates on a data structure that satisfies the key requirement for safe exception propagation: it is a directed acyclic graph (DAG) rooted at the human principal, in which every edge is immutable and every node carries a human anchor identifier. The protocol traverses this DAG upward — never modifying it, never relying on runtime process state, and never depending on a machine actor's goodwill to preserve the chain's integrity.

The Zylos Research analysis of multi-agent orchestration observes that hierarchical decomposition enables localized replanning — when a sub-agent's plan fails, only the subtree containing the failure need be replanned, not the entire task hierarchy \cite{zylos-htn}. This property is orthogonally complementary to the Bailout Protocol: localized replanning governs how tasks are retried and re-delegated under normal execution; the Bailout Protocol governs what happens when a node encounters an exception that exceeds its replanning authority. The two mechanisms are non-overlapping and jointly exhaustive: the agent system handles expected failures through its planning and replanning mechanisms; the Bailout Protocol handles unexpected failures through mandatory escalation to a human principal. Together they cover the full space of failure modes without relying on a machine actor to decide which category a failure falls into.

\end{document}
